\section{Stationary profiles and the contradiction}\label{sec:stationary}

We retain the bad-event laws \(Q_N\) and the adapted episodes from
Section~\ref{sec:stages}. Their entropy and drop bounds
\eqref{eq:34}, \eqref{eq:38}, and \eqref{eq:40} will produce a stationary
sequence of local endpoint profiles with positive mean normalization
loss. We begin by identifying the statistic that will make this
stationary law impossible.

For a full-support subprobability \(w\) on horizontal space, define
\[
 T_R(w)=\max_{j\in\mathbb Z}
       \min\{w(z_2\le j),w(z_2\ge j+R)\},\qquad R\in\mathbb Z_{\ge1}.
\]
Here \(z_2\) denotes the second full-space coordinate. This positive,
translation-invariant quantity measures how much mass can be placed in
each of two tails separated by \(R\). The profile update multiplies
surviving mass by the episode normalization factors. The outward-order
estimate~\eqref{eq:33} transports the two tails at a cost logarithmic
in the traversed height, together with a local environmental remainder.
For a stationary marked profile law, we will derive
\[
 mE c_0\le A_\to E\log\left(1+\sum_{i=0}^{m-1}h_i\right)+O(1),
\]
with the last term independent of \(m\). The logarithmic height moment
from the episode bounds makes the right side \(o(m)\), contradicting
\(E c_0>0\).

The construction must supply a profile to which this argument applies.
Local limits can initially have diffuse mass or several components
escaping from one another. The array below retains those possibilities,
then finite-window entropy excludes diffuse mass and reduces the
positive-loss event to finitely many persistent components of full
support. Marking one component uniformly preserves stationarity. The
same entropy bound controls its upper environment relative to an iid
field while retaining exactly its current profile and offset data.

Translation-invariant compactifications that retain separated components
and diffuse mass were developed by Mukherjee and
Varadhan~\cite[Theorem~3.2]{MukherjeeVaradhan2016} for occupation measures
and by Bates and Chatterjee~\cite[Theorem~2.8]{BatesChatterjee2020} for
polymer endpoint distributions. We construct the array needed here
directly, storing sampled-anchor offsets together with episode marks
and the corresponding upper environments.

To remove the environmental bias from the local remainder, we first fix
a finite number \(m\) of episodes and send the tail separation \(R\)
to infinity. Only afterwards do we send \(m\) to infinity. No mixing
estimate uniform in the episode window is needed. The final profile
argument establishes the displayed obstruction in this order.

\subsection{The stationary array}


Horizontal space is \(\mathbb Z^{d-1}\). We observe profiles from sampled
anchors, allowing offsets between anchors to escape to infinity in the
limit.

Pad each episode sequence to all integer indices by \(t_i=0\) for \(i<0\), \(t_i=N\) for \(i\ge M\), and zero marks \(h_i,J_i,c_i\) outside \(0\le i<M\). Write \(\mu_i^*=\mu_{t_i}\) for the normalized global endpoint profile there, viewed as a probability on horizontal sites. Conditional on the entire environment, sample anchors \(U_{i,a}\sim\mu_i^*\) independently for indices \(i\in\mathbb Z\) and \(a\ge1\). Keep the following array in a countable compact product:
\begin{itemize}
\item The nonnegative integer marks \(h_i,J_i\) and nonnegative real marks
\(c_i\), each range compactified by \(+\infty\).
\item The offsets \(O_{ia,jb}=U_{j,b}-U_{i,a}\), in the lattice with its
one-point compactification, and all profile coordinates
\(w_j^{ia}(z)=\mu_j^*(U_{i,a}+z)\).
\item The upper environments \(\Omega^{i,m,a}(l,z)\), for \(m\ge1\),
\(l\ge0\), and horizontal \(z\). For \(l<t_{i+m}-t_i\), use the true row
\(\omega((t_i+l,U_{i,a}+z),\cdot)\). At the remaining heights use the
row at \((l,U_{i,a}+z)\) of a full auxiliary iid upper field for the
pair \((i,m)\). This field is common to all anchors \(a\), and the
auxiliary fields are independent across pairs and independent of the
environment and anchors.
\end{itemize}
Denote by \(P_\uparrow\) the iid upper-field law on indices \((l,z)\),
including \(l=0\). The functions \(q(0;\Omega)\) and
\(a_H(0;\Omega)\) are the earlier quenched no-drop functions evaluated
in this upper field; they require no rows below its base.

Profile and row coordinates use the ordinary compact ranges \([0,1]\) and the row simplex, respectively. Let \(\tau\) shift the array by reading all absolute episode indices one index later (leaving relative offsets in time such as \(m\) unchanged). Take the probability law \(\mathbb L_N\) averaging shifts to episode starts:
\[
\mathbb L_N[F]=(E_{Q_N}M)^{-1} E_{Q_N,{\rm aux}}\sum_{i=0}^{M-1}F(\tau^i{\rm array}) .
\]
\begin{lemma}[Stationary episode limit]\label{profile-stationary}
Along a subsequence, \(\mathbb L_N\) converges weakly to a shift-invariant
law \(\mathbb L\). Under this law, simultaneously for every \(i\),
\[
\begin{gathered}
1\le h_i<\infty,\qquad J_i,c_i<\infty,\\
E_{\mathbb L}\bigl(c_i+\log(1+h_i)\bigr)<\infty,
\qquad E_{\mathbb L}c_i>0.
\end{gathered}
\]
Every stored row is uniformly elliptic with the original bound
\(\kappa\).
\end{lemma}

\begin{proof}
Compactness gives a subsequential weak limit. The shift is continuous,
and the occupation sum telescopes: for every bounded continuous \(F\),
\[
\left|\mathbb L_N[F\circ\tau]-\mathbb L_N[F]\right|
\le \frac{2\|F\|_\infty}{E_{Q_N}M}\longrightarrow0.
\]
At the origin, \eqref{eq:40} bounds the occupation mean of \(J_i\).
Truncation passes this bound to the limit, so \(J_0<\infty\) almost
surely. The bounds \eqref{eq:38} then pass and give the stated
integrability and finiteness. For the positive truncated-drop bound in
\eqref{eq:40}, put $C_{L_0}=B_0/k+2bL_0$ and use the globally
bounded continuous test
\[
 (c_0\wedge C_{L_0})\mathbf1_{\{J_0\le L_0\}}.
\]
It agrees with the required drop on every supported array, by
\eqref{eq:38}. Its expectation therefore passes to the limit and gives
$E_{\mathbb L}c_0>0$. The occupation law gives \(h_0\ge1\),
and shift invariance gives all assertions at every integer index.
Uniform ellipticity is a closed condition on the row coordinates and
therefore also passes to the limit.
\end{proof}

Write
\[
H_{i,m}=\sum_{s=i}^{i+m-1}h_s,\qquad C_{i,m}=\sum_{s=i}^{i+m-1}c_s .
\]
By the stationary ergodic theorem, in its possibly nonergodic form
\cite[Theorem~6.2.1]{Durrett2019}, there exists \(\zeta>0\) and a
shift-invariant positive-probability event
\[
\mathcal E=\{\lim_{m\to\infty}C_{0,m}/m>\zeta\}
\tag{41}\label{eq:41}
\]
(up to null sets).

\subsection{Entropy in a finite episode window}

We use standard relative-entropy tools originating with Kullback and
Leibler~\cite{KullbackLeibler1951}; the binary-event and projection
inequalities are instances of data processing, as discussed in
Sason and Verd\'u~\cite[Section~II-B]{SasonVerdu2016}. Their needed forms,
including the variational formula, are proved below before being applied
to the stopped episode windows.

Current data for entropy purposes are only
\[
\mathcal A_i=\big((O_{ia,ib})_{a,b},(w_i^{ia}(z))_{a,z}\big).
\]
The reference law keeps the marginal of \(\mathcal A_i\).
Conditionally on these data, labels at finite offsets share one iid
upper field, translated by their offsets, while distinct classes use
independent iid fields. The finite-offset relations in the limit are
consistent equivalence relations with additive offsets: each specified
finite-offset event is clopen, so every such consistency identity passes
to the limit. Write \(\mathbf R(\mathcal A_i,d\phi)\) for this reference
conditional field law.

\begin{lemma}[Finite-window entropy]\label{profile-entropy}
There is a constant \(C_E<\infty\) such that, for every fixed \(i\) and
\(m\ge1\), the \(\mathbb L\)-law of
\((\mathcal A_i,(\Omega^{i,m,a})_a)\) has relative entropy at most
\(C_E m\) with respect to its own current-data marginal followed by
the reference kernel \(\mathbf R\).
\end{lemma}

\begin{proof}
We first record the elementary entropy facts used here. Relative
entropy is \({\rm KL}(Q\mid P')=\int f_0\log f_0\,dP'\), where
\(f_0=dQ/dP'\), and is infinite without absolute continuity. It has
the variational formula
\[
{\rm KL}(Q\mid P')
=\sup_{F\text{ bounded measurable}}
\left\{QF-\log P'e^F\right\}.
\]
The upper bound follows from Jensen applied to \(e^F/f_0\).
For equality, truncate \(\log f_0\) between \(\pm B\) and send
\(B\to\infty\), or test a singular set. Thus projection cannot
increase entropy. On compact metric spaces, bounded continuous tests
suffice by bounded approximation in \(L^1(Q+P')\); the variational
formula therefore gives lower semicontinuity under joint weak
convergence of both laws.

For an event of respective probabilities \(s,t\), where \(0<t<1\),
projection to its indicator gives
\({\rm KL}(Q\mid P')\ge s\log(1/t)-\log2\).
If a reference law is a data marginal times an independent fixed law,
replacing that marginal by the controlled law's \emph{own} data
marginal only decreases finite entropy. Indeed the new density is
the old density divided by its conditional mean given the data, and
the entropy difference is the nonnegative marginal entropy. These
log-density identities are taken on sets of positive controlled
density; negative parts are integrable, and conditional Jensen for
\(f_0\log f_0\) justifies the marginal term.

Before the limit, at deterministic \(i\ge0\), put
\(\mathcal H_i=\sigma(\mathcal F_{t_i},(U_{i,a})_a)\) and
\(\Phi_i=(\Omega^{i,m,a})_a\); \(\mathcal A_i\) is
\(\mathcal H_i\)-measurable. Include the same auxiliary sampling under
\(P\). Given \(\mathcal H_i\), the spliced field above \(t_i\), in
absolute horizontal coordinates, is still iid under \(P\). Indeed, at
both bounded layer stopping times the layers at and above the stopping
layer are fresh, while the current anchors use only lower-layer data.
Keeping the slab up to the second stopping layer and replacing its tail
by an independent copy therefore preserves the iid law. Thus \(\Phi_i\)
has conditional law \(\mathbf R(\mathcal A_i,d\phi)\).

Let \(L_t=dQ_N|_{\mathcal F_t}/dP|_{\mathcal F_t}\). Under \(Q_N\),
the conditional density of the information stopped at \(t_{i+m}\),
together with the independent extension, relative to the corresponding
\(P\)-law given \(\mathcal H_i\), is \(L_{t_{i+m}}/L_{t_i}\).
The anchors at \(i\) use the same lower-data sampling kernel under both
laws, and the extension costs no density. Conditional Jensen therefore
gives, for every bounded \(F\),
\[
\begin{aligned}
&E_{Q_N,{\rm aux}}[F(\mathcal A_i,\Phi_i)\mid\mathcal H_i]\\
&\quad\le
E_{Q_N,{\rm aux}}[\log L_{t_{i+m}}-\log L_{t_i}\mid\mathcal H_i]
+\log\int e^{F(\mathcal A_i,\phi)}\mathbf R(\mathcal A_i,d\phi).
\end{aligned}
\]
Multiply by \(\mathbf1_{\{i<M\}}\), which is past-known, and average as for \(\mathbb L_N\). Jensen to pull the log outside this average gives the \(\mathbb L_N\) entropy bound at index zero (against its own current-data marginal and the kernel) of
\[
\frac{1}{E_{Q_N}M}\sum_{i\ge0}E_{Q_N}\big[\log L_{t_{i+m}}-\log L_{t_i}\big]
\le \frac{mD_0\log N}{E_{Q_N}M}.
\]
Increments from \(i\ge M\) are zero. The displayed sum telescopes
over deterministic indices, since \(t_i=N\) for \(i\ge N\); its
subtracted initial expected log densities are nonnegative relative
entropies. For passage to the limit, finite-class consistency is a
closed condition: self-offsets vanish, finite offsets are symmetric
with sign change, and specified finite offsets add on triples.
Violations are detected by finitely many clopen coordinate events.
On this support the reference kernel is weakly continuous. For any
finitely many field coordinates, coincidence of the input sites is
specified by a finite-offset equality, while different sites have
independent rows. Finite-coordinate tests suffice by compactness.
Consequently the reference laws converge weakly with their current-data
marginals. Lower semicontinuity and \eqref{eq:40} yield the bound
\(C_E m\); shift invariance gives every \(i\).
\end{proof}

\subsection{Profiles and multiplicities}

Offsets across all times determine equivalence classes of labels
\((i,a)\) by finite distance. Within one class, the profiles at any
time \(j\), viewed from its different anchors, are translates of one
another and have total mass at most one. These consistency statements
pass coordinatewise.

\begin{lemma}[Sampling identities]\label{profile-sampling}
Under \(\mathbb L\), simultaneously for all indices and horizontal
\(z\),
\[
2^{-n}\sum_{b\le2^n}\mathbf1_{\{O_{ia,jb}=z\}}\longrightarrow w_j^{ia}(z)\quad\text{a.s.}
\tag{42}\label{eq:42}
\]
A distinct label \((j,b)\ne(i,a)\) cannot have finite offset \(z\)
from \((i,a)\) while \(w_j^{ia}(z)=0\).
\end{lemma}

\begin{proof}
Before the limit, the anchor draws at \(j\), excluding a possible
own label, are conditionally independent samples given the environment
and \(U_{i,a}\). The occupation indices and weights depend only on
the environment, not on these draws. The mean-square error in
\eqref{eq:42} is therefore \(O(2^{-n})\), including the possible
own-label contribution. The finite-offset indicators are continuous,
so these summable bounds pass to the limit. Borel--Cantelli proves
\eqref{eq:42}.

For distinct labels and any continuous function \(g\) on \([0,1]\),
the same conditional sampling identity gives
\[
E_{\mathbb L}\!\left[
g(w_j^{ia}(z))
\bigl(\mathbf1_{\{O_{ia,jb}=z\}}-w_j^{ia}(z)\bigr)\right]=0.
\]
Approximate the indicator of zero by bounded continuous functions
to obtain the second assertion. All the indices and sites are
countable, so the assertions hold on one common probability-one event.
\end{proof}

\begin{lemma}[Profile updates and averaged survival]\label{profile-updates}
For every \(i,m,a\), write \(H=H_{i,m}\), \(C=C_{i,m}\), and
\(\Omega=\Omega^{i,m,a}\). Coordinatewise,
\[
w_{i+m}^{ia}\ \ge\ e^C\, w_i^{ia} K_H^\Omega .
\tag{43}\label{eq:43}
\]
Here the kernel maps horizontal starting positions at layer zero to
endpoints at layer \(H\) in the anchor frame. Also, almost surely,
\[
\limsup_n 2^{-n}\sum_{a\le2^n} a_{H_{i,m}}(0;\Omega^{i,m,a})\le e^{-C_{i,m}}.
\tag{44}\label{eq:44}
\]
\end{lemma}

\begin{proof}
Before the limit, local no-drop continuation is contained in the global
survival stream. Dividing its endpoint inequality by the new global
mass gives \eqref{eq:43}, with
\(C=\log(a_{t_i}/a_{t_{i+m}})\). At finite limiting marks, pass first
the inequalities involving finitely many starting sites and
bounded-length successful kernel paths. Strict violations of these
finite-path inequalities are open, since each path weight is a finite
product of row coordinates. Increasing the finite restrictions gives
\eqref{eq:43}. Successful paths to layer \(H\) use no departure row
at or above \(H\), so the auxiliary tail is irrelevant.

Before the limit, at a fixed index the average in \eqref{eq:44} has
conditional mean at most \(e^{-C}\) given the environment, again by
the global mass inequality. Its terms are independent bounded
functions of the independently sampled current anchors; they do not
use appended rows. Chebyshev bounds the conditional probability of
exceeding \(e^{-C}+\epsilon\) by \(2^{-n}\epsilon^{-2}\), and this
also holds after occupation averaging. For fixed finite \(H\), the
function \(a_H\) is lower semicontinuous, being a sum of nonnegative
continuous finite-path weights. Thus the strict upper-violation event
is open at finite marks, and the same summable probability bound
passes to the limit. Borel--Cantelli, followed by a countable sequence
\(\epsilon\downarrow0\), proves \eqref{eq:44}.
\end{proof}

Call a current label \emph{proper} if \(w_i^{ia}(0)>0\), and
\emph{dust} otherwise.

\begin{proposition}[Finitely many persistent profiles]\label{profile-classes}
On \(\mathcal E\), there is no dust. The proper classes form the same
finite nonempty set at every time, and the profile of each class has
full support on horizontal space.
\end{proposition}

\begin{proof}
\emph{Dust is isolated.}
If a dust label \((i,a)\) had a distinct finite-offset neighbor
\((j,b)\), with offset \(z\), profile consistency would give
\(w_i^{jb}(-z)=w_i^{ia}(0)=0\). Apply the second assertion of
Lemma~\ref{profile-sampling} with \((j,b)\) as the reference label
and \((i,a)\) as the sampled target. This is impossible. Thus dust
labels are singleton classes even when all times are included.

At any fixed time, the dust indicators are exchangeable, because
anchor-permutation symmetry passes to the limit. They have an
asymptotic frequency \(\delta_i\), which is positive if any dust occurs. To see
the latter assertion, apply the ergodic theorem to the stationary
label sequence of dust indicators. On the invariant zero-frequency
event, taking expectations shows that each indicator vanishes.

Fix \(i,m\). Under the reference law of Lemma~\ref{profile-entropy},
conditional on \(\mathcal A_i\), the singleton dust classes have
independent upper fields. Their values \(q(0;\Omega^{i,m,a})\) are
therefore iid with mean \(p\). Finite entropy gives absolute
continuity, so the strong law on the dust subsequence also holds under
\(\mathbb L\) whenever that subsequence is infinite. Since
\(a_H\ge q\), including in every spliced field, \eqref{eq:44} yields
\[
e^{-C_{i,m}}\ge p\delta_i.
\]
This holds simultaneously for every \(m\). On \(\mathcal E\), its
left side tends to zero, by \eqref{eq:41} and stationarity. Dust is
therefore absent at all times on \(\mathcal E\).

\emph{There are finitely many proper classes.}
For fixed \(n\) and \(r_0>0\), whenever possible choose from
\(\mathcal A_i\) labels in \(n\) distinct classes with self masses
at least \(r_0\). On this selection event, when
\(C_{i,m}>\zeta m\), \eqref{eq:43} and the total-mass bound for
each output profile give
\[
\prod_{\rm chosen} q(0;\Omega^{i,m,a})\le r_0^{-n}e^{-n\zeta m}.
\]
The event of successful selection together with this deterministic
product bound is measurable from the current data and the extension
fields. Its reference probability is at most
\[
r_0^{-n\lambda}(E q^{-\lambda})^n e^{-\lambda n\zeta m},
\]
by \eqref{eq:4} and class independence. The binary entropy bound and
Lemma~\ref{profile-entropy} show that its \(\mathbb L\)-probability
has \(\limsup_{m\to\infty}\) at most \(C_E/(\lambda n\zeta)\).
On the selection event intersected with \(\mathcal E\), the product
bound holds for all sufficiently large \(m\). This bounds the
probability of that intersection by the same quantity. Send
\(r_0\downarrow0\), then \(n\to\infty\), to exclude infinitely
many proper classes.

\emph{The classes persist and have full support.}
Over any positive integer height, the local kernel has positive weight
between any two horizontal sites: move horizontally at the starting
layer and then ascend directly. Uniform ellipticity makes this finite
script positive. Thus \eqref{eq:43} gives a full-support next profile
for every current proper class. By \eqref{eq:42}, every positive
coordinate is represented by next-time labels. Distinct classes cannot
merge, by the global finite-offset equivalence relation.

Conditional on the shift-invariant event \(\mathcal E\), the current
number of classes is therefore nondecreasing, finite, positive, and
stationary. It is consequently constant almost surely. Indeed, writing $N_i$ for
this count, stationarity and monotonicity give
\[
 E[\min(N_{i+1},K)-\min(N_i,K)]=0
\]
for every integer $K$. The nonnegative difference vanishes almost
surely; letting $K$ increase proves $N_{i+1}=N_i$. This uses no moment
of the number of classes. The class sets transport bijectively
through time, and every profile has full support by the update from
the preceding time.
\end{proof}

\subsection{Why a stationary profile cannot sustain normalization growth}

The preceding construction has left us with a finite, persistent set of
full-support profiles on the event $\mathcal E$. Condition on
$\mathcal E$, mark one class uniformly, and use its least current label
as the origin at each episode. Uniform
marking commutes with time shift, so the marked law $\mathbb L_*$ is
stationary. This choice matters: a class sampled according to its mass
would generally receive a different weight after an episode.

For any fixed window of $m$ episodes, put
\[
 H_m=\sum_{i=0}^{m-1}h_i,\qquad C_m=\sum_{i=0}^{m-1}c_i.
\]
Write $w_0$ for the marked profile at the start of this window and write
$w_m$ in the same horizontal frame. Both are subprobability measures of
full support. The profile update gives
\[
 w_m\ge e^{C_m}w_0 K_{H_m}^{\Omega_m}
 \quad\hbox{coordinatewise},
\]
where $\Omega_m$ is the spliced upper field viewed from the current
representative. The profile at time $m$, with its usual representative,
has the same distribution as $w_0$ up to translation.

There is one further property of this marked field that we will need.
Let $D_*=(\mathcal A_0,a)$, where $\mathcal A_0$ is exactly the
current offset and profile data defined above and $a$ is the least
current label of the marked class. In particular, $D_*$ contains no
future episode marks or future profiles, and $w_0$ is a function of
$D_*$. The joint law of $(D_*,\Omega_m)$ has finite relative entropy with respect to
\begin{equation*}\tag{45}\label{eq:45}
 \mathbb L_*(D_*)\otimes P_\uparrow.
\end{equation*}
Indeed, conditioning on $\mathcal E$ changes the density of the projected
unmarked law by a factor at most $1/\mathbb L(\mathcal E)$. Finite entropy
therefore remains finite. Given the current data, uniform selection
among the current class representatives is a probability kernel. Apply
that same kernel to the reference law in Lemma~\ref{profile-entropy};
the selected field has iid law $P_\uparrow$ independently of the selected
index and the current data. Projection and the entropy chain rule now
give \eqref{eq:45} with the controlled law's own data
marginal.

The following argument isolates the remaining obstruction. It requires
finite entropy separately for each fixed $m$; no entropy bound uniform
in $m$ is used.

\begin{proposition}[No positive stationary normalization rate]
\label{profile-no-growth}
Suppose a stationary marked profile law has full-support subprobability
profiles and nonnegative marks $c_i$, positive integer marks $h_i$, with
\[
 E[c_0+\log(1+h_0)]<\infty.
\]
Assume the coordinatewise update \eqref{eq:43}, the finite
entropy property \eqref{eq:45} for every fixed $m$, and
the outward-order estimate \eqref{eq:33}, including its uniform local
approximation property. Then $Ec_0=0$.
\end{proposition}

\begin{proof}
We use the two-tail statistic introduced at the start of this section.
For an integer $R\ge1$, recall that
\[
 T_R(w)=\max_{j\in\mathbb Z}
 \min\{w(z_2\le j),w(z_2\ge j+R)\}.
\]
Here $z_2$ means the second full-space coordinate of a horizontal site.
Full support and finite mass imply $0<T_R(w)\le1$, and a maximizing
integer $j_R$ exists. Choose one measurably. For each fixed full-support
$w$, its maximizing cuts satisfy
\[
 j_R\longrightarrow-\infty,\qquad j_R+R\longrightarrow+\infty.
\]
For example, if $j_R$ were bounded below along a subsequence, the left
mass at both $j_R$ and $j_R-1$ would stay bounded below whereas the right
mass would tend to zero. Moving the cut one unit left would strictly
increase the smaller mass, contradicting maximality. The other endpoint
is treated by moving the cut right. Thus the normalized restrictions
$\alpha_-$ and $\alpha_+$ of $w_0$ to these two tails eventually sample
sites outside every fixed neighborhood of the representative origin.

Fix $m$ throughout the next steps. Let $Q_R$ be the mass of endpoint
pairs $(u,v)$ with $v_2-u_2\ge R$ under
\[
 (\alpha_-K_{H_m}^{\Omega_m})\otimes
 (\alpha_+K_{H_m}^{\Omega_m}).
\]
Every starting pair already has signed gap at least $R$. Its outward
order event is therefore contained in the event just tested. Finite
entropy transfers the probability-one event in \eqref{eq:33} to
$\Omega_m$, simultaneously for the countably many starting pairs and
integer heights, so it applies at the random height $H_m$. That estimate
and Jensen's inequality yield
\begin{equation*}\tag{46}\label{eq:46}
\begin{aligned}
 \log Q_R&\ge-A_\to\log(1+H_m)-Z_R,\\
 Z_R&=\iint\mathcal Z((0,x),(0,y);\Omega_m)
                \,\alpha_-(dx)\alpha_+(dy).
\end{aligned}
\end{equation*}
Write $\nu_\pm=\alpha_\pm K_{H_m}^{\Omega_m}$. There is an integer
$j$ with $\nu_-(u_2\le j)\ge Q_R/2$ and
$\nu_+(v_2\ge j+R)\ge Q_R/2$. Choose $j$ minimal for the first
inequality, so $\nu_-(u_2<j)<Q_R/2$. If the second inequality failed,
split ordered pairs into $u_2<j$ and $u_2\ge j$. The latter condition
forces $v_2\ge j+R$. Each part would then have mass strictly less than
$Q_R/2$, since both endpoint measures have mass at most one, contrary
to the definition of $Q_R$.

Each initial tail has mass at least $T_R(w_0)$. Applying
\eqref{eq:43} to the two tails at this cut gives
\begin{equation*}\tag{47}\label{eq:47}
 \log T_R(w_m)-\log T_R(w_0)
 \ge C_m-A_\to\log(1+H_m)-Z_R-\log2.
\end{equation*}
This is the decisive inequality: normalization contributes $C_m$,
whereas preserving the order of the two tails costs only a logarithm
of the traversed height and the local environmental remainder.

We next remove the environmental bias from that remainder. We claim
\begin{equation*}\tag{48}\label{eq:48}
 \limsup_{R\to\infty}E Z_R
 \le M_\mathcal Z:=\sup_{x,y}E_P\mathcal Z(x,y)<\infty.
\end{equation*}
Write $f_m(D,\Omega)$ for the density relative to
\eqref{eq:45}. Under that product reference, additionally
sample $x,y$ from $\alpha_-\otimes\alpha_+$ conditionally on $D$ alone.
Then $EZ_R$ is the reference expectation of
$f_m(D,\Omega)\mathcal Z((0,x),(0,y);\Omega)$.

The exponential moment in \eqref{eq:33} is uniform in the sampled pair.
Consequently the density and the remainder can both be truncated with
errors uniform in $R$. To see the only delicate part, on $\{f_m>B\}$
Young's inequality gives, for sufficiently small $\lambda>0$,
\[
 \lambda f_m\mathcal Z\le f_m\log f_m+e^{\lambda\mathcal Z}.
\]
The entropy term has vanishing tail. The reference probability of
$\{f_m>B\}$ is at most $1/B$, and Cauchy--Schwarz bounds the exponential
term on that event by $O(B^{-1/2})$. Once $f_m$ is bounded, truncating
$\mathcal Z$ is justified directly by its exponential moment.

Approximate the bounded density by a bounded nonnegative function
$g(D,\Omega)$ using only a fixed finite box of field rows around the
origin, allowing arbitrary dependence on $D$. With $L^1$ error at most
$\varepsilon$, its reference integral is at most $1+\varepsilon$.
Approximate the capped remainder, uniformly in $x,y$, by a bounded
function using fixed-radius
upper-row neighborhoods of those two sites. This follows from the
local approximation in \eqref{eq:33}; boundedness converts convergence
in probability to convergence in $L^1$. As $R$ increases, the two tail
supports and their neighborhoods leave the fixed origin box. Under the
product reference, conditional on $D,x,y$, the two approximating
factors then use disjoint iid row sets. Their product expectation is
bounded by $(M_\mathcal Z+\varepsilon)E_{P_\uparrow}g(D,\cdot)$.
The threshold in $R$ may depend on $D$; boundedness permits integration
by reverse Fatou. Sending the approximation errors to zero and removing
the truncations proves \eqref{eq:48}. This entire step
keeps $m$ fixed, so its approximation boxes need not be uniform in $m$.

\begin{figure}[tbp]
\centering
\begin{tikzpicture}[x=1cm,y=1cm,font=\small,>=Stealth]
  \fill[blue!7] (-5.65,0) rectangle (-3.65,2.15);
  \fill[green!7] (3.65,0) rectangle (5.65,2.15);
  \fill[black!8] (-1.35,0) rectangle (1.35,1.65);
  \draw[blue!60!black,densely dashed] (-5.65,0) rectangle (-3.65,2.15);
  \draw[green!40!black,densely dashed] (3.65,0) rectangle (5.65,2.15);
  \draw[black!55,densely dashed] (-1.35,0) rectangle (1.35,1.65);
  \node[align=center,text=blue!60!black] at (-4.65,1.12)
    {local rows\\near $(0,x)$};
  \node[align=center] at (0,.9) {fixed origin box\\for $g(D_*,\Omega)$};
  \node[align=center,text=green!40!black] at (4.65,1.12)
    {local rows\\near $(0,y)$};
  \draw[->,black!65] (-6.1,0) -- (6.1,0) node[right] {$x_2$};
  \fill (0,0) circle (2pt);
  \node[anchor=north,align=center] at (0,-.08) {representative\\origin};
  \fill[blue!60!black] (-4.65,0) circle (2.5pt);
  \node[anchor=north,text=blue!60!black] at (-4.65,-.08) {$x\sim\alpha_-$};
  \fill[green!40!black] (4.65,0) circle (2.5pt);
  \node[anchor=north,text=green!40!black] at (4.65,-.08) {$y\sim\alpha_+$};
  \draw[black!55] (-3.1,-.08) -- (-3.1,.12);
  \node[anchor=north] at (-3.1,-.08) {$j_R$};
  \draw[black!55] (3.1,-.08) -- (3.1,.12);
  \node[anchor=north] at (3.1,-.08) {$j_R+R$};
  \draw[<-,blue!60!black,line width=1pt] (-6,-.85) -- (-3.1,-.85);
  \node[anchor=north,text=blue!60!black] at (-4.65,-.95)
    {$x_2\le j_R\longrightarrow-\infty$};
  \draw[->,green!40!black,line width=1pt] (3.1,-.85) -- (6,-.85);
  \node[anchor=north,text=green!40!black] at (4.65,-.95)
    {$y_2\ge j_R+R\longrightarrow+\infty$};
  \draw[<->,black!55] (-5.65,2.6) -- (5.65,2.6);
  \node[anchor=south,align=center] at (0,2.7)
    {$\mathcal Z_{\rm loc}(x,y;\Omega)$ uses only the two outer neighborhoods};
\end{tikzpicture}

\caption{The finite-row approximations in the proof of
Equation~\eqref{eq:48}. At a fixed episode range, the
bounded density approximation uses a fixed box around the representative
origin. The local remainder uses neighborhoods of the two sampled tail
sites. These neighborhoods leave the origin box as the tail parameter
increases, so the approximating factors are conditionally independent
under the product reference law. Only the second coordinate is shown.}
\label{fig:tail-decorrelation}
\end{figure}

For fixed $R,m$, the right side of \eqref{eq:47} is
integrable. The entropy inequality and the same exponential moment
show that $EZ_R<\infty$. Let
$X=\log T_R(w_0)$ and $Y=\log T_R(w_m)$. Translation invariance of $T_R$
and stationarity make these finite nonpositive variables identically
distributed. Their difference has integrable negative part by
\eqref{eq:47}. Clipping both variables below at $-B$ gives
an integrable zero-mean difference; its positive and negative parts
increase to the corresponding parts of $Y-X$. Monotone convergence
therefore shows that $Y-X$ is integrable and $E(Y-X)=0$. Taking
expectations in \eqref{eq:47}, and now sending $R$ to
infinity, yields
\[
 mEc_0\le A_\to E\log(1+H_m)+M_\mathcal Z+\log2.
\]
Finally let $m$ increase. With $Z_i=\log(1+h_i)$,
\[
 \log(1+H_m)\le\log m+\max_{0\le i<m}Z_i,
 \qquad
 E\max_{i<m}Z_i\le T+mE[Z_0\mathbf1_{\{Z_0>T\}}].
\]
Divide by $m$, send $m$ to infinity and then $T$ to infinity. The assumed
logarithmic moment gives $E\log(1+H_m)=o(m)$, hence $Ec_0\le0$.
Nonnegativity of $c_0$ proves the proposition.
\end{proof}

\begin{proof}[Completion of the proof of Theorem~\ref{thm:main}]
Apply Proposition~\ref{profile-no-growth} to $\mathbb L_*$. The profile
construction supplies its update and moment hypotheses, and the marked
field argument above supplies the finite entropy hypothesis. But on the
shift-invariant event $\mathcal E$, the stationary ergodic theorem gives
$E_{\mathbb L_*}c_0>\zeta$. This contradicts the proposition and excludes
the bad-event sequence $Q_N$. Thus \eqref{eq:5} holds.
Proposition~\ref{geom-speed} gives finite mean regeneration duration
and an almost-sure deterministic vector velocity with positive first
coordinate. This vector lies on the deterministic ray of
Corollary~\ref{geom-ray}, whose projection on the original transient
direction $\ell$ is positive. Undoing the signed permutation of the
coordinate axes therefore gives $X_n/n\to v$ with $v\cdot\ell>0$,
as asserted in Theorem~\ref{thm:main}.
\end{proof}
